\documentclass[aspectratio=169,10pt]{beamer}
\usepackage{default_latex_beamer_template}

\title[NY Reduced Model Progress]{NY Reduced Power-System Model:\\Current Progress}
\subtitle{Validation of a 49-bus research model for interface-flow and preliminary dynamic line-rating studies}
\author{NPCC / NY Reduced-Model Workstream}
\date{Status through July 14, 2026 \quad $\vert$ \quad Prepared July 19, 2026}

\begin{document}

\defaulttitleframe{A self-contained update for technical audiences new to the project}

\begin{frame}{Executive Summary}
\small
\begin{columns}[T,onlytextwidth]
\begin{column}{0.64\linewidth}
\begin{itemize}
  \item A \pass{49-bus New York reduced model is now accepted} for the summer-peak operating point and two reactive-load sensitivity tests.
  \item The accepted model passes AC power flow, generator and voltage limits, trusted branch limits, interface-flow accuracy, dispatch-movement limits, and a full-to-reduced structural benchmark.
  \item Four transfer interfaces are calibrated, and five physical transmission circuits have current results that can support later conductor mapping.
  \item \fail{Winter peak and high-Total-East conditions remain excluded} because they exceed fixed dispatch and/or interface-accuracy limits.
\end{itemize}
\end{column}
\begin{column}{0.33\linewidth}
\metric{green}{ACCEPTED}{Summer-peak research case}
\vspace{0.35em}
\metric{blue}{71.7 MW}{Calibrated-interface MAE}
\vspace{0.35em}
\metric{amber}{5 circuits}{Physical current scope}
\end{column}
\end{columns}

\begin{alertblock}{Claim boundary}
{\scriptsize The model supports preliminary dynamic line-rating methodology research. It does not yet contain the physical and operational data needed for utility planning or real-time operating claims.}
\end{alertblock}
\vspace{-0.24cm}
\source{Current validation, interface, dispatch, and structural benchmark ledgers}
\end{frame}

\begin{frame}{What Are We Building?}
\begin{tikzpicture}[
  node distance=0.58cm,
  box/.style={rounded corners=2pt,minimum height=1.65cm,text width=3.65cm,align=center,inner sep=0.38cm,draw=blue,fill=paleblue},
  arrow/.style={-{Stealth[length=2.5mm]},very thick,blue}
]
\node[box] (source) {\textbf{Regional source model}\par\scriptsize Detailed topology, generators, controls, and operating data};
\node[box,right=of source,fill=palegreen,draw=green] (reduced) {\textbf{NY-focused reduced model}\par\scriptsize 49 retained buses, boundary equivalents, explicit interfaces};
\node[box,right=of reduced,fill=paleamber,draw=amber] (research) {\textbf{Research outputs}\par\scriptsize Voltages, transfers, circuit currents, and future weather-based ratings};
\draw[arrow] (source) -- (reduced);
\draw[arrow] (reduced) -- (research);
\end{tikzpicture}

\vspace{0.55cm}
\begin{columns}[T,onlytextwidth]
\begin{column}{0.49\linewidth}
\begin{block}{Reduced model}
A smaller electrical equivalent that preserves the behavior needed for selected New York studies while removing most of the surrounding regional detail.
\end{block}
\end{column}
\begin{column}{0.49\linewidth}
\begin{block}{Dynamic line rating}
A weather-dependent current limit. This deck validates electrical currents; conductor and weather inputs are a later layer.
\end{block}
\end{column}
\end{columns}
\source{Current model architecture and validation documentation}
\end{frame}

\begin{frame}{How the Project Moved Forward}
\centering
\begin{tikzpicture}[scale=0.84,transform shape,
  node distance=0.45cm,
  box/.style={rounded corners=2pt,minimum height=1.55cm,text width=3.25cm,align=left,inner sep=0.40cm},
  arrow/.style={-{Stealth[length=2.5mm]},very thick,blue}
]
\node[box,fill=palered,draw=red] (old) {\textbf{Starting blocker}\par\small Public-load cases did not produce certified feasible solutions.};
\node[box,fill=paleblue,draw=blue,right=of old] (controls) {\textbf{Controls repaired}\par\small Source-backed status, voltage targets, reactive limits, and reference control were preserved.};
\node[box,fill=paleamber,draw=amber,right=of controls] (network) {\textbf{Reduction improved}\par\small Explicit interfaces and physical circuits were added; non-physical equivalents were corrected.};
\node[box,fill=palegreen,draw=green,right=of network] (now) {\textbf{Current result}\par\small Summer peak passes every mandatory gate; two broader conditions remain excluded.};
\draw[arrow] (old) -- (controls);
\draw[arrow] (controls) -- (network);
\draw[arrow] (network) -- (now);
\end{tikzpicture}

\vspace{0.65cm}
\begin{block}{The major process change}
Cases are no longer accepted because they merely converge. Publication now requires electrical feasibility, realistic control provenance, bounded redispatch, interface accuracy, and full-to-reduced structural agreement.
\end{block}
\source{Project handoff and current validation policy}
\end{frame}

\begin{frame}{Current Model Configuration}
\small
\begin{columns}[T,onlytextwidth]
\begin{column}{0.49\linewidth}
\begin{block}{Network and operating point}
\begin{itemize}
  \item 49 buses, 83 branches, 49 generator records
  \item New York zones A--K plus external boundary equivalents
  \item Five explicitly mapped physical circuits
  \item Summer load: 30,644.99 MW
  \item Generation: 31,433.46 MW
  \item Active loss: 788.47 MW
\end{itemize}
\end{block}
\end{column}
\begin{column}{0.49\linewidth}
\begin{block}{Controls and measurements}
\begin{itemize}
  \item 35 regulating generators at seven eligible buses
  \item Source-backed active and reactive limits
  \item Fixed external active-power schedules
  \item Fixed-zero external reactive schedule by default
  \item Four calibrated interface operators
  \item Three additional interface proxies retained as diagnostics
\end{itemize}
\end{block}
\end{column}
\end{columns}

\begin{exampleblock}{Why the circuit count increased}
Two retained corridors are represented as separate parallel circuits, producing 83 model branches rather than 81 while preserving five one-to-one physical circuit rows.
\end{exampleblock}
\source{Current case dimensions, control map, and dispatch/loss ledger}
\end{frame}

\begin{frame}{What the Validation Gates Mean}
\scriptsize
\begin{tabularx}{\linewidth}{>{\bfseries}p{2.75cm} X >{\raggedleft\arraybackslash}p{2.5cm} >{\raggedleft\arraybackslash}p{2.0cm}}
\toprule
Gate family & Requirement & Current result & Status \\
\midrule
Power flow & Standard and reactive-limit AC power flow & both converge & \pass{PASS} \\
Restoration & No artificial P/Q/voltage/branch slack & 0 & \pass{PASS} \\
Physical limits & Voltage, generator P/Q, trusted branch violations & 0 / 0 / 0 / 0 & \pass{PASS} \\
Power balance & Reference adjustment below 1 MW & $1.17\times10^{-6}$ MW & \pass{PASS} \\
Interfaces & MAE $\leq100$; max $\leq200$; $|$bias$|\leq50$ MW & 71.7 / 98.2 / 28.8 & \pass{PASS} \\
Dispatch & Movement $\leq10\%$ and $\leq3000$ MW & 6.433\%; 1,971 MW & \pass{PASS} \\
Controls & No unexpected regulating buses; reference matched & 0; matched & \pass{PASS} \\
Structural match & Current, loss, and net-Q errors each $\leq15\%$ & all below 15\% & \pass{PASS} \\
\bottomrule
\end{tabularx}

\vspace{0.55em}
\begin{block}{Margin indicators}
Trusted static maximum loading is 56.3\% of its stated rating, and no regulating group is at a reactive limit. Voltages reach approximately 0.90--1.10 pu, so hard bounds pass but preferred voltage margin remains narrow.
\end{block}
\source{Current mandatory and preferred validation ledger}
\end{frame}

\begin{frame}{Accepted Summer-Peak Interface Results}
\begin{columns}[T,onlytextwidth]
\begin{column}{0.62\linewidth}
\begin{tikzpicture}
\begin{axis}[
  xbar,
  width=0.94\linewidth,
  height=5.45cm,
  xmin=-120,xmax=120,
  axis x line*=bottom,
  axis y line*=left,
  symbolic y coords={Moses South,West Central,Central East,Dysinger East},
  ytick=data,
  xlabel={Target error (MW)},
  grid=major,
  grid style={draw=muted!18},
  tick label style={font=\scriptsize},
  label style={font=\scriptsize},
  nodes near coords,
  nodes near coords align={horizontal},
  every node near coord/.append style={font=\scriptsize},
  enlarge y limits=0.18
]
\addplot[fill=teal,draw=teal] coordinates {
  (-85.74,Moses South)
  (65.37,West Central)
  (37.35,Central East)
  (98.19,Dysinger East)
};
\end{axis}
\end{tikzpicture}
\end{column}
\begin{column}{0.35\linewidth}
\metric{blue}{71.7 MW}{Mean absolute error}
\vspace{0.42em}
\metric{green}{98.2 MW}{Maximum error}
\vspace{0.42em}
\metric{teal}{28.8 MW}{Signed bias}
\vspace{0.55em}
{\scriptsize All aggregate gates pass: MAE $\leq100$, maximum $\leq200$, and $|$bias$|\leq50$ MW.}
\end{column}
\end{columns}
\source{Current summer-peak interface validation ledger; calibrated operators only}
\end{frame}

\begin{frame}{Operating Conditions Tested}
\scriptsize
\begin{tabularx}{\linewidth}{X c r r r r l}
\toprule
Condition & Reactive load & IF MAE & Max IF & Dispatch & Load & Decision \\
 & & MW & MW & MW & \% & \\
\midrule
Summer peak & Nominal & 71.7 & 98.2 & 1,971 & 6.433 & \pass{Accepted} \\
Summer peak & 90\% & 71.5 & 98.1 & 1,971 & 6.433 & \pass{Accepted} \\
Summer peak & 110\% & 71.6 & 98.1 & 1,971 & 6.433 & \pass{Accepted} \\
\addlinespace
Winter peak & All tests & 41.4 & 80.9 & 3,199 & 13.602 & \fail{Excluded} \\
High Total East & All tests & 156.8 & 318.4 & 3,410 & 12.037 & \fail{Excluded} \\
\bottomrule
\end{tabularx}

\vspace{0.65cm}
\begin{columns}[T,onlytextwidth]
\begin{column}{0.49\linewidth}
\begin{alertblock}{Winter peak}
Electrical and interface tests pass, but the case requires more than 3,000 MW of zonal redispatch and more than 10\% of total load.
\end{alertblock}
\end{column}
\begin{column}{0.49\linewidth}
\begin{alertblock}{High Total East}
The case exceeds both redispatch limits and the calibrated-interface accuracy limits.
\end{alertblock}
\end{column}
\end{columns}
\source{Current dispatch, loss, and interface validation ledgers}
\end{frame}

\begin{frame}{Full-to-Reduced Structural Benchmark}
\begin{columns}[T,onlytextwidth]
\begin{column}{0.67\linewidth}
\begin{tikzpicture}
\begin{axis}[
  ybar,
  width=\linewidth,
  height=5.3cm,
  ymin=0,ymax=17,
  bar width=15pt,
  xtick={1,2,3,4,5},
  xticklabels={GL,PV--WS 1,PV--WS 2,WS--MW 1,WS--MW 2},
  ylabel={Circuit-current error (\%)},
  grid=major,
  grid style={draw=muted!18},
  tick label style={font=\scriptsize},
  label style={font=\scriptsize},
  nodes near coords,
  every node near coord/.append style={font=\scriptsize},
]
\addplot[fill=blue,draw=blue] coordinates {(1,5.051) (2,7.555) (3,7.555) (4,11.302) (5,11.302)};
\draw[red,dashed,very thick] (axis cs:0.5,15) -- (axis cs:5.5,15);
\node[anchor=south east,font=\scriptsize,text=red] at (axis cs:5.5,15) {15\% gate};
\end{axis}
\end{tikzpicture}
\end{column}
\begin{column}{0.30\linewidth}
\metric{green}{0.000096\%}{Active-loss error}
\vspace{0.40em}
\metric{blue}{2.470\%}{Net-network-Q error}
\vspace{0.40em}
\metric{teal}{PASS}{Reactive-limit power flow}
\end{column}
\end{columns}

\begin{block}{What this demonstrates}
For the same 2019 operating snapshot, the reduced network matches selected physical-circuit currents, active losses, and net reactive behavior within the adopted 15\% research threshold.
\end{block}
\source{2019 full-to-reduced structural benchmark ledger}
\end{frame}

\begin{frame}{Dynamic-Line-Rating Circuit Scope}
\begin{columns}[T,onlytextwidth]
\begin{column}{0.58\linewidth}
\scriptsize
\begin{tabularx}{\linewidth}{X c r}
\toprule
Physical corridor & Circuit count & Summer current \\
\midrule
Gilboa--Leeds & 1 & 1.157 kA \\
Pleasant Valley--Wood Street & 2 & 0.925 kA each \\
Wood Street--Millwood & 2 & 0.591 kA each \\
\bottomrule
\end{tabularx}

\vspace{0.65em}
\begin{exampleblock}{What is usable now}
The five mapped rows carry terminal and series currents with physical-circuit provenance suitable for later conductor assignment.
\end{exampleblock}
\end{column}
\begin{column}{0.39\linewidth}
\begin{alertblock}{What is not yet included}
\begin{itemize}
  \item Conductor type and ampacity
  \item Thermal parameters and weather
  \item Route segmentation and sag/clearance
  \item Terminal-equipment ratings
  \item Dynamic ampacity calculation
\end{itemize}
\end{alertblock}
\end{column}
\end{columns}

\begin{block}{Interpretation rule}
Currents on aggregate and boundary equivalents are electrical diagnostics only. A model branch rating is not automatically a conductor ampacity or dynamic rating.
\end{block}
\source{Current branch-current and physical-circuit mapping ledgers}
\end{frame}

\begin{frame}{Remaining Issues and Limitations}
\small
\begin{columns}[T,onlytextwidth]
\begin{column}{0.49\linewidth}
\begin{alertblock}{Operating coverage}
\begin{itemize}
  \item Winter peak and high-Total-East conditions fail fixed acceptance gates.
  \item Total East, UPNY--ConEd, and Dunwoodie South remain diagnostic because their complete monitored-circuit definitions are not conserved.
  \item Summer voltages reach approximately 0.90/1.10 pu: hard bounds pass, but preferred margin is narrow.
\end{itemize}
\end{alertblock}
\end{column}
\begin{column}{0.49\linewidth}
\begin{alertblock}{Reduction and physical fidelity}
\begin{itemize}
  \item The non-physical equivalent correction is empirical rather than an exact Ward reduction.
  \item Three control locations remain provisional fixed-reactive-power buses.
  \item One legacy equivalent rating overload remains diagnostic because its rating meaning is not trusted.
  \item No conductor or weather layer has been attached.
\end{itemize}
\end{alertblock}
\end{column}
\end{columns}

\begin{block}{Claim boundary}
The model is a defensible preliminary research equivalent, not a utility planning, reliability, facility-rating, real-time dispatch, or operating model.
\end{block}
\source{Current limitations and mandatory-gate documentation}
\end{frame}

\begin{frame}{Recommended Next Work}
\begin{enumerate}
  \item \textbf{Complete interface definitions:} replace diagnostic proxies with explicit conserved circuit or cutset operators where the retained topology supports them.
  \item \textbf{Broaden operating coverage without relaxing gates:} reduce winter and high-Total-East redispatch, and improve high-Total-East interface closure.
  \item \textbf{Increase voltage-control margin:} investigate buses near 0.90/1.10 pu and resolve the three provisional control locations without placeholder reactive limits.
  \item \textbf{Add the physical rating layer:} map conductors, equipment limits, route segments, weather, thermal models, and sag/clearance constraints to the five circuits.
  \item \textbf{Revalidate before publication:} rerun all operating and structural gates and publish only cases that pass every mandatory test.
\end{enumerate}

\begin{exampleblock}{Near-term deliverable}
An updated validation package that preserves the accepted summer-peak model and either passes the broader operating conditions under unchanged gates or keeps them explicitly excluded.
\end{exampleblock}
\source{Current documented limitations and acceptance policy}
\end{frame}

\begin{frame}[plain]
\begin{tikzpicture}[remember picture,overlay]
  \fill[navy] (current page.south west) rectangle (current page.north east);
  \fill[teal] (current page.south west) rectangle ([xshift=0.18cm]current page.north west);
\end{tikzpicture}
\vspace{1.18cm}
{\color{white}\Huge\bfseries Current takeaway\par}
\vspace{0.65cm}
{\color{paleblue}\Large Use the accepted summer-peak model for preliminary dynamic line-rating methodology development.\par}
\vspace{0.45cm}
{\color{white}\large Keep winter peak and high-Total-East conditions excluded. Do not make utility-operational or dynamic-rating claims until interface coverage and physical conductor/weather inputs are complete.\par}
\vfill
{\color{paleblue}\small Status through July 14, 2026}
\end{frame}

\end{document}
